\chapter{Conclusion}
\label{chap:conclusion}

This thesis asked whether a patient-specific liver digital twin, built from
clinical data and graded intervention evidence, can support the
evidence-grounded evaluation and ranking of existing recipes for an individual
patient (Section~\ref{sec:research_question}). The answer is yes, within the
limits set out below: HepatoTwin estimates a patient's liver state from
routine data with calibrated probabilities, links that state to graded
evidence, and ranks real recipes so that the top of every list meets the
patient's nutritional references, cuisine and practical limits far more often
than the alternatives. This chapter summarises what each objective achieved,
states the limitations, and sets out future work.

% ======================================================================
\section{Summary of Findings}
\label{sec:c_findings}

\paragraph{Liver state from routine data (Objective 1).} The liver-state model
detects steatosis with an AUC of 0.881 for CAP $\geq$ 248~dB/m and 0.879 for
CAP $\geq$ 280~dB/m on held-out participants, with an expected calibration
error of at most 0.022, and outperforms the Fatty Liver Index and the Hepatic
Steatosis Index on the same patients. Its 90\% conformal interval covered
90.9\% of test participants. On NHANES~III, collected in 1988--1994 and
assessed by ultrasound, it reached an AUC of 0.771, and the quarter of
participants with the highest predicted probability had 1.58 times the
all-cause and 2.15 times the cardiovascular mortality. The point estimate of
CAP, in contrast, explains only 56\% of its variance (Chapter~\ref{chap:cap-experiments}),
so the twin shows patients calibrated chances rather than a single number.

\paragraph{Evidence-graded priorities (Objective 2).} Every priority cites
rows of a table of 25 findings, each linking an intervention to an outcome
with an explicit strength label, and a causal knowledge graph of 16 links, and none carries an invented
numeric target. The patient's own 24-hour recall moves a priority by at most
one level, which keeps a noisy single recall from dominating.

\paragraph{Recipe evaluation and ranking (Objective 3).} The health-gated
ranking placed recipes that meet every nutritional reference, the patient's
cuisine and the practical limits in 82.3\% of top-ten lists, against 27.2\%
for weighted scoring without the gate, and every recommended recipe met every
health reference. Personal fit raised the number of distinct top-ten lists
from 28 to 124 while that share barely changed (82.3\% against 82.6\%), and
every seven-day plan was
complete and met every constraint.

\paragraph{Scenarios and updating (Objective 4).} In a simulation calibrated
to published weight-loss trials, the twin's weight forecasts erred by 1.9~kg
three months ahead and 2.8~kg six months ahead, against 2.9 and 4.7~kg for a
standard textbook rule, with 91\% of outcomes inside its stated ranges. When
it named ``plan not followed'' as the cause of slow progress it was always
correct, and it tracked liver fat with
an AUC of 0.88 at months 0, 6 and 12.

\paragraph{Honest evaluation (Objective 5).} The evaluation also exposed what
the twin cannot do, and the system was changed in response. A patient's body
response to diet cannot be told apart from ordinary weight drift within a
year (correlation 0.03--0.06), so the twin reports a lack of progress rather
than its cause and asks a clinician to review. The model is less well
calibrated for some ethnic groups; FIB-4 is weak in young community patients,
so a measured liver stiffness is preferred when available; and the first
plain-language summaries were faithful but hard to read, which led to a
rewrite that lowered the reading grade from 12.5 to 9.8 while keeping every
statement faithful to the model output.

% ======================================================================
\section{Limitations}
\label{sec:c_limitations}

\begin{itemize}
    \item \textbf{No real follow-up data.} NHANES examines each participant
          once. The current-state estimate is validated on held-out patients,
          but the twin's projections and its updating were evaluated only on
          virtual patients. They are labelled as simulated projections, not
          clinical forecasts, and remain unvalidated against repeat
          measurements in real patients.
    \item \textbf{A ceiling on prediction from routine data.} Liver fat
          measured by CAP is only partly explained by routine measurements;
          no model tried exceeded an $R^2$ of about 0.56
          (Chapter~\ref{chap:cap-experiments}).
    \item \textbf{Calibration differs between groups.} Discrimination was good
          in every group (AUC 0.84--0.93), but the predicted probability was
          too high for non-Hispanic Black participants and, at the borderline
          of significance, too low for non-Hispanic Asian participants.
    \item \textbf{Evidence is at the level of diets, not recipes.} The
          evidence table describes the effects of weight loss and dietary
          patterns; no single recipe can claim those effects, which is why the
          twin ranks recipes against reference values rather than predicting
          their effect on the liver.
    \item \textbf{Recipe data.} RecipeDB reports total rather than free sugars,
          some diet labels do not match the ingredients (so recipes must pass
          both the label and an ingredient check), and main meals are unevenly
          spread across cuisines, which limits weekly plans for some cuisines.
    \item \textbf{Narrow external validation.} The NHANES~III analysis is
          unweighted and adjusted for age and sex only, because smoking,
          alcohol and income are not mapped.
    \item \textbf{No user study.} The planned pilot evaluation with patients
          and clinicians was not carried out within the project's time, so the
          usability and acceptability of the dashboard and its recommendations
          have not been tested with real users.
\end{itemize}

% ======================================================================
\section{Future Work}
\label{sec:c_future}

\begin{itemize}
    \item \textbf{A pilot with patients and clinicians}, to test whether the
          recommendations are acceptable and cookable, whether the summaries
          and the liver visualisation are understood, and whether clinicians
          find the patient page useful when approving plans.
    \item \textbf{Longitudinal deployment}, in which patients enter their own
          measurements over months, with repeat FibroScan examinations, so that
          the twin's projections and its re-assessment loop can be validated on
          real data.
    \item \textbf{Recalibration by group}, refitting the calibrators where the
          fairness audit found the probabilities too high or too low.
    \item \textbf{A broader recipe corpus}, with more main meals in
          under-represented cuisines, and recipe modification by ingredient
          substitution, so that a familiar dish can be made to meet the
          references rather than replaced.
    \item \textbf{Objective adherence data}, such as step counts from a phone
          or wearable, in place of self-reported walking at check-ins.
\end{itemize}

% ======================================================================
\section{Closing Remarks}
\label{sec:c_closing}

HepatoTwin shows that a digital twin for MASLD can be built from data
collected in a routine check-up and linked, step by step, to the evidence
behind each dietary recommendation. Its estimates are calibrated, its
recommendations are traceable to graded evidence, and its limits are stated
rather than hidden. Validating its projections on real patients over time is
the step that would turn it from a research prototype into a tool for care.
